These supplementary analyses provide operational and configuration-sensitivity context. The scheduling simulation is separate from the classification experiments and does not measure cloud-QPU queue latency. The configuration diagnostics use Breast Cancer as specified in their captions and are exploratory; matched end-to-end resource measurements would be needed to use them in a deployment ranking.

\subsection{Scheduling Sensitivity}
\label{sec:scheduling}

Table~\ref{tab:scheduling} reports a separate 50-job scheduling simulation comparing FIFO, shortest-job-first (SJF), and priority scheduling across one, two, and four QPUs. The reported waiting times and makespans provide supplementary operational context for workload scheduling. These results do not represent measured cloud-QPU queue latency or classification accuracy.

\begin{table}[H]
  \centering
  \caption{Supplementary scheduling comparison (50 jobs, varying QPU count)}
  \label{tab:scheduling}
  \small
  \setlength{\tabcolsep}{4pt}
  \resizebox{\columnwidth}{!}{%
    \begin{tabular}{cccc}
      \toprule
      \textbf{Strategy} & \textbf{QPUs} & \textbf{Avg Wait (ms)} & \textbf{Makespan (ms)} \\
      \midrule
      FIFO     & 1 & $1156 \pm 87$ & 2464 \\
      SJF      & 1 & $\phantom{0}508 \pm 63$ & 2464 \\
      Priority & 1 & $1143 \pm 92$ & 2464 \\
      \midrule
      FIFO     & 2 & $457 \pm 41$ & 1251 \\
      SJF      & 2 & $183 \pm 29$ & 1317 \\
      Priority & 2 & $462 \pm 45$ & 1275 \\
      \midrule
      FIFO     & 4 & $119 \pm 18$ & 666 \\
      SJF      & 4 & $\phantom{00}54 \pm 11$ & 756 \\
      Priority & 4 & $142 \pm 21$ & 690 \\
      \bottomrule
    \end{tabular}
  }
  \begin{tablenotes}
    \small
    \item Values reported as mean $\pm$ std over 10 scheduling simulations.
  \end{tablenotes}
\end{table}

\FloatBarrier

\subsection{Configuration Diagnostics}
\label{sec:ablation}

The following one-factor comparisons use Breast Cancer with an 8-qubit circuit. The default uses angle encoding, $L=2$, Adam, and $\epsilon_{2q}=10^{-2}$. These comparisons are exploratory. The scripts and execution logs that produced Tables~\ref{tab:ablation_encoding}, \ref{tab:ablation_optimizer}, \ref{tab:ablation_pca}, and~\ref{tab:ablation_loss} are not available in the retained materials, so it cannot be independently confirmed that these diagnostics used the final circuit architecture of Figure~\ref{fig:vqc_ansatz}. Table~\ref{tab:ablation_depth} is analytical and is computed from the final architecture. The other diagnostics also differ from the primary protocol. They use Adam or SGD with parameter-shift gradients (or COBYLA) for roughly 60--85 iterations and PCA retaining 70\% of the variance, and the number of trials behind their $\pm$ values is not recorded. Their accuracies (about 0.72--0.79 for the VQC and 0.91--0.92 for XGBoost) are therefore not comparable with the primary 20-iteration SPSA results in Table~\ref{tab:vqc_accuracy} (0.4474 for this configuration) or with the matched classical results in Table~\ref{tab:matched_comparison}. The higher VQC accuracies are consistent with the primary SPSA budget being training-limited (Section~\ref{sec:experiments}). Because their provenance cannot be verified, however, they are not used as evidence for that interpretation. Raw trial data, leakage controls, and configuration metadata will support future verification. The supplementary analyses identify further inputs for a deployment decision.

\subsubsection{Encoding and Ansatz Depth}

Table~\ref{tab:ablation_encoding} reports encoding summaries. Physical state-preparation complexity provides the basis for an encoding-depth and hardware-feasibility comparison. Table~\ref{tab:ablation_depth} gives quantities derived from the stated ansatz and Eq.~\eqref{eq:fidelity}. Trial-level metadata can support further depth-accuracy analysis.

\begin{table}[H]
  \centering
  \caption{Supplementary encoding diagnostic (Breast Cancer, $n=8$, $L=2$, Adam)}
  \label{tab:ablation_encoding}
  \small
  \setlength{\tabcolsep}{4pt}
  \resizebox{\columnwidth}{!}{%
    \begin{tabular}{lccc}
      \toprule
      \textbf{Encoding} & \textbf{Acc ($\epsilon_{2q}=10^{-3}$)} & \textbf{Acc ($\epsilon_{2q}=10^{-2}$)} & \textbf{Encode Gates} \\
      \midrule
      Angle     & $0.764 \pm 0.033$ & $0.768 \pm 0.037$ & $n$ (1Q only) \\
      Amplitude & $0.778 \pm 0.035$ & $0.726 \pm 0.040$ & $O(2^n)$ \\
      \bottomrule
    \end{tabular}
  }
\end{table}
\FloatBarrier

\begin{table}[H]
  \centering
  \caption{Supplementary ansatz-depth diagnostic (Breast Cancer, $n=8$, angle, Adam, $\epsilon_{2q}=10^{-2}$)}
  \label{tab:ablation_depth}
  \small
  \setlength{\tabcolsep}{4pt}
  \resizebox{\columnwidth}{!}{%
    \begin{tabular}{lcc}
      \toprule
      \textbf{Layers $L$} & \textbf{2Q Gates} & \textbf{Analytical gate-survival proxy (\%)} \\
      \midrule
      1 & \phantom{0}7 & 90.27 \\
      2 & 14 & 82.80 \\
      3 & 21 & 75.95 \\
      4 & 28 & 69.67 \\
      5 & 35 & 63.90 \\
      \bottomrule
    \end{tabular}
  }
\end{table}
\FloatBarrier

\subsubsection{Optimizer, PCA, and Loss Diagnostics}

Tables~\ref{tab:ablation_optimizer}, \ref{tab:ablation_pca}, and~\ref{tab:ablation_loss} give optimizer, preprocessing, and loss sensitivity summaries. The MSE and cross-entropy rows in Table~\ref{tab:ablation_loss} are supplementary Adam diagnostics, not the primary SPSA results. Dedicated factorial experiments can estimate factor importance. A parameter-shift evaluation count gives a circuit-evaluation accounting quantity. Provider-job counts require provider records. Fold-local fitting documentation and historical loss verification can strengthen the preprocessing and loss analyses.

\begin{table}[H]
  \centering
  \caption{Supplementary optimizer diagnostic (Breast Cancer, $n=8$, $L=2$, angle, $\epsilon_{2q}=10^{-2}$)}
  \label{tab:ablation_optimizer}
  \small
  \setlength{\tabcolsep}{4pt}
  \resizebox{\columnwidth}{!}{%
    \begin{tabular}{lccc}
      \toprule
      \textbf{Optimizer} & \textbf{Accuracy} & \textbf{Iters to Conv.} & \textbf{Circuit evaluations} \\
      \midrule
      SGD    & $0.740 \pm 0.042$ & $\sim$85 & $2PB\,N_{\text{iter}}$ \\
      Adam   & $0.768 \pm 0.037$ & $\sim$62 & $2PB\,N_{\text{iter}}$ \\
      COBYLA & $0.722 \pm 0.045$ & $\sim$71 & $B\,N_{\text{iter}}$ \\
      \bottomrule
    \end{tabular}
  }
  \begin{tablenotes}
    \small
    \item $P=n(2L+1)=40$ trainable parameters and batch size $B$, so parameter-shift uses $2PB$ circuit evaluations per gradient step; COBYLA is gradient-free and uses $B$ evaluations (one loss evaluation) per iteration. Iterations to convergence are approximate.
  \end{tablenotes}
\end{table}
\FloatBarrier

\begin{table}[H]
  \centering
  \caption{Supplementary PCA/feature-selection diagnostic (Breast Cancer, $n=8$, $L=2$, Adam, $\epsilon_{2q}=10^{-2}$)}
  \label{tab:ablation_pca}
  \small
  \setlength{\tabcolsep}{4pt}
  \resizebox{\columnwidth}{!}{%
    \begin{tabular}{lccc}
      \toprule
      \textbf{Feature Strategy} & \textbf{VQC Accuracy} & \textbf{XGBoost Accuracy} & \textbf{Gap} \\
      \midrule
      Full PCA (70\% var.) & $0.768 \pm 0.037$ & $0.908 \pm 0.010$ & 0.140 \\
      MI Feature Selection & $0.789 \pm 0.040$ & $0.917 \pm 0.010$ & 0.128 \\
      \bottomrule
    \end{tabular}
  }
  \begin{tablenotes}
    \small
    \item Point-estimate gaps are descriptive; preprocessing is insufficiently specified for inferential use.
  \end{tablenotes}
\end{table}
\FloatBarrier

\begin{table}[H]
  \centering
  \caption{Supplementary loss-function diagnostic (Breast Cancer, $n=8$, $L=2$, Adam, $\epsilon_{2q}=10^{-2}$)}
  \label{tab:ablation_loss}
  \small
  \setlength{\tabcolsep}{4pt}
  \resizebox{\columnwidth}{!}{%
    \begin{tabular}{lcccc}
      \toprule
      \textbf{Loss} & \textbf{Accuracy} & \textbf{Iters to Conv.} & \textbf{Extra QPU Cost} & \textbf{Gap to XGBoost} \\
      \midrule
      MSE & $0.768 \pm 0.037$ & $\sim$62 & None (same $\langle Z_0\rangle$ estimates) & 0.140 \\
      Cross-entropy & $0.781 \pm 0.035$ & $\sim$58 & None (same $\langle Z_0\rangle$ estimates) & 0.127 \\
      \bottomrule
    \end{tabular}
  }
  \begin{tablenotes}
    \small
    \item Actual additional evaluations, shots, and job submissions must be reported before comparing loss-function efficiency.
  \end{tablenotes}
\end{table}
\FloatBarrier

Within these legacy diagnostics, the point estimates differ by 0.046 across optimizers and by 0.042 across encodings at $\epsilon_{2q}=10^{-2}$, and by 0.013--0.021 across loss function and feature strategy. The primary 15-split analysis (Table~\ref{tab:noise_sensitivity}) finds that $\epsilon_{2q}$ changes mean accuracy by at most 0.006 for a fixed model and 0.061 with retraining; the two analyses use different optimizers and are not directly comparable. Dedicated causal studies can estimate factor importance. The primary deployment conclusions use the matched primary evidence.